\scanpage{174}
\mspage{48}

On a donc
\diagram{\begin{tikzcd}[ampersand replacement=\&,row sep=1.2em,column sep=2em]
 \underset{\alpha\in J}{\underrightarrow{\mathrm{lim}}}\,Z_\alpha \arrow[r] \arrow[d,phantom,"\unsure{\wr|{\cong}}" description] \&
 \underset{\alpha\in\mathrm{Ss}_{\unsure{\pi|\varpi}}X}{\underrightarrow{\mathrm{lim}}}\,Z_\alpha \arrow[r,"\varphi"] \& X,\\
 X \& \&
\end{tikzcd}}
le composé étant l'identité. Si on sait
que $\varphi$ est un mono, il en résulte
que c'est un \unsure{iso} (mono. ayant
une section). \deleted{\unsure{Or}} Or $\varphi$ \unsure{sera} un
mono si \unsure{tte|toute|\deleted{tout}} $\underrightarrow{\mathrm{lim}}$\unsure{;|:} filtrante de
mono. est mono, \inserted{et alors on gagne.}
\deleted{puisque $\varphi$ est la limite inductive des
$(Z_\alpha\xrightarrow{i_\alpha} X'_\alpha\inserted{=X})$ correspondant à l'hom\unsure{omorphisme}
évident du syst. inductif $(Z_\alpha)_{\alpha\in\mathrm{Ss}_{\pi}X}$ dans le
\unsure{const}.}

Si on suppose \deleted{\unsure{donc}} que les
éléments de $\mathrm{Filt}^{\pi}$ \unsure{ont|sont} $\pi$-\unsure{accessible|accessibles},
on gagne aussi, car on trouve

\uline{Corollaire 5.4.} Si \deleted{\unsure{les}} objets de $\pi$
sont $\pi$-accessibles (p. ex. si $\pi\geq\pi_1$ et $\pi=\pi^{\pi_0}$,
p. ex. $\pi\geq\pi_1$ et $\pi=2^c$, avec $c\geq\pi_0$, cf
(3.2.) p. 16), alors pour tte partie filtrante
$J$ de $\mathrm{Ss}(X)$ grande devant $\pi$, telle que
$\underset{\alpha\in J}{\unsure{\mathrm{lim}|\underrightarrow{\mathrm{lim}}}}\,Z_\alpha\overset{\sim}{\longrightarrow}X$,
$J$ est cofinale dans $\mathrm{Ss}_{\pi}X$,
(donc $\underset{\alpha\in J}{\underrightarrow{\mathrm{lim}}}\,Z_\alpha\overset{\sim}{\longrightarrow}
\underset{\alpha\in\mathrm{Ss}_{\pi}(X)}{\underrightarrow{\mathrm{lim}}}\,Z_\alpha$)\deleted{\unsure{)}}.
\deleted{En effet, \unsure{soit}} C'est immédiat.

% archive sheet number: 173
